\documentclass[11pt,a4paper]{article}
\usepackage[utf8]{inputenc}
\usepackage[T1]{fontenc}
\usepackage[spanish,es-tabla]{babel}
\usepackage[margin=2.5cm]{geometry}
\usepackage{lmodern}
\usepackage{booktabs}
\usepackage{longtable}
\usepackage{tabularx}
\usepackage{array}
\usepackage{enumitem}
\usepackage[hidelinks]{hyperref}
\setlength{\parskip}{0.5em}
\setlength{\parindent}{0pt}
\newcolumntype{Y}{>{\raggedright\arraybackslash}X}

\title{Recuperaci\'on de im\'agenes por caracter\'isticas:\\ planeaci\'on del proyecto}
\author{Juan Ulises Calderon Huerta}
\date{12 de octubre de 2026}

\begin{document}
\maketitle

\section{Objetivo y alcance}

Comparar dos estrategias para describir objetos en im\'agenes y recuperarlos por sus caracter\'isticas: (A) \textbf{atributos fijos}, predichos por clasificadores entrenados, y (B) \textbf{descripci\'on generada} por un modelo visi\'on--lenguaje. Dataset principal: \textbf{VeRi-776}. Fecha l\'imite: \textbf{9 de noviembre de 2026}, con avances cada lunes.

\textbf{Incluido:}
\begin{itemize}[noitemsep]
  \item Enfoque A: clasificador multicabeza para color y tipo de veh\'iculo.
  \item Enfoque B: modelo generativo preentrenado, sin entrenamiento.
  \item Localizaci\'on con YOLO26, evaluada contra las cajas anotadas de VeRi.
  \item Prueba secundaria con LocateAnything (\texttt{q8\_0}) en GPU.
  \item Conjunto sint\'etico de prueba con fondos y posiciones variados.
   \item Base de datos \textbf{MongoDB} (probable) para guardar los resultados y consultar por caracter\'isticas (por ejemplo, ``SUV rojo'').
  \item Comparaci\'on de calidad, latencia y memoria en el equipo de pruebas.
\end{itemize}

\textbf{Fuera de alcance:} ajuste fino de modelos generativos, detector entrenado desde cero, b\'usqueda con embeddings y despliegue. DeepFashion queda condicionado al hito del 26 de octubre.

\section{Equipo de pruebas}

\begin{tabular}{@{}ll@{}}
\toprule
CPU & Intel Core i7-10750H @ 2.60\,GHz \\
GPU & NVIDIA GeForce RTX 2070 Max-Q (8\,GB VRAM) \\
RAM & 16\,GB \\
Almacenamiento & 1\,TB NVMe PCIe 3.0 \\
\bottomrule
\end{tabular}
\newpage

\section{Modelos y Herramientas}

Criterio: que cada experimento quepa en 8\,GB de VRAM y se pueda repetir en este equipo.

\begin{table}[h]
\centering
\small
\begin{tabularx}{\textwidth}{@{}lYY@{}}
\toprule
Modelo & Papel & Nota \\
\midrule
YOLO26n, YOLO26s & Localizaci\'on principal & Versi\'on m\'as reciente de YOLO; variantes ligeras. \\
EfficientNet-B0 & Atributos fijos (principal) & Unos 5.3\,M de par\'ametros. \\
MobileNetV3-Small & Atributos fijos (ligero) & Unos 2.5\,M de par\'ametros. \\
Qwen2.5-VL-3B-Instruct & Descripci\'on generada (principal) & Se probar\'a cuantizado a 4 bits. \\
Florence-2-base & Descripci\'on generada (alternativa) & Se usa si Qwen no funciona en una jornada. \\
LocateAnything-3B (\texttt{q8\_0}) & Localizaci\'on secundaria & Archivo de unos 6.3\,GB; se ejecuta por separado de los dem\'as modelos. \\
MongoDB  & Base de datos de resultados & Un documento por objeto detectado: imagen, caja, atributos, descripci\'on y modelo usado. \\
\bottomrule
\end{tabularx}
\caption{Modelos a utilizar.}
\end{table}


\section{Actividades y experimentos}

Los m\'etodos y configuraciones pueden ajustarse durante el desarrollo.

\begin{itemize}[noitemsep]
  \item \textbf{Datos:} auditar VeRi, usar la partici\'on oficial por identidad y revisar condiciones de uso (uso no comercial).
  \item \textbf{Base de datos y recuperaci\'on:} guardar los resultados en MongoDB (un documento por objeto) y consultar con atributos predichos por cada enfoque.
  \item \textbf{Referencia:} recuperaci\'on con etiquetas reales.
  \item \textbf{Atributos fijos:} EfficientNet-B0 (solo cabezas y ajuste limitado) y MobileNetV3-Small.
  \item \textbf{Descripci\'on generada:} Qwen2.5-VL-3B cuantizado, o Florence-2-base como alternativa.
  \item \textbf{Localizaci\'on:} YOLO26n y YOLO26s preentrenados; ajuste fino con cajas de VeRi si hay margen de mejora.
  \item \textbf{Datos sint\'eticos:} recortes de veh\'iculos sobre fondos variados y en distintas posiciones, para evaluar robustez.
  \item \textbf{LocateAnything:} pruebas con distintos \emph{prompts} para verificar si localiza y entrega descriptores espec\'ificos.
  \item \textbf{Base de datos y recuperaci\'on:} guardar los resultados en MongoDB (un documento por objeto) y consultar con atributos predichos por cada enfoque.
  \item \textbf{Comparaci\'on:} calidad, latencia y memoria en el equipo descrito.
  \item \textbf{Opcional:} DeepFashion, ResNet18.
\end{itemize}

\newpage

\section{Cronograma}\label{sec:cronograma}

Entregas los lunes 12, 19 y 26 de octubre, y 2 y 9 de noviembre.

\begin{longtable}{@{}p{3cm}p{7cm}p{4.5cm}@{}}
\caption{Cronograma.}\\
\toprule
Periodo & Actividades & Entrega \\
\midrule
\endfirsthead
\toprule
Periodo & Actividades & Entrega \\
\midrule
\endhead
12 oct. & Alcance, plan y selecci\'on de modelos. & Planeaci\'on. \\
\midrule
13--18 oct. & Auditar VeRi y preparar recortes. YOLO26 sin ajuste. Primer modelo generativo. Preparar pruebas de LocateAnything. & 19 oct.: auditor\'ia de datos y primeros resultados. \\
\midrule
20--25 oct. & Referencia con etiquetas reales. Atributos fijos. Modelo generativo. Ajuste de YOLO26 si aplica. Generador sint\'etico. & 26 oct.: primera comparaci\'on de calidad, latencia y memoria. \\
\midrule
26 oct. & Decidir si se incluye DeepFashion. & Decisi\'on registrada. \\
\midrule
27 oct.--1 nov. & Base de datos y consultas. Pipeline con detecci\'on autom\'atica. Prueba de robustez con datos sint\'eticos. LocateAnything. Congelar modelos el 1 nov. & 2 nov.: demostraci\'on funcional. \\
\midrule
3--8 nov. & Evaluaci\'on final, informe, demostraci\'on y repositorio. & 9 nov.: entrega final. \\
\bottomrule
\end{longtable}

\textbf{Si hay retrasos}, se recorta en este orden: DeepFashion, ResNet18 y MobileNetV3-Small, datos sint\'eticos, ajuste de YOLO26 y LocateAnything.

\end{document}